\chapter{项目概览与阅读指南}

\section{先用一句话理解项目}

ECom Review Intelligence Agent 是一个面向 B 端电商运营的评论证据决策原型。你可以把它理解成一位“会查评论、会整理证据、但不会替运营拍板”的助手：运营提出一个商品问题，系统返回可回查的评论证据，帮助单品诊断、商品比较和问题分流。它不是面向 C 端消费者的购物推荐助手，也不会根据评论自动推断真实缺陷率、销量或利润。

\begin{beginnerbox}
先记住项目的完整故事：P0、P1 先把数据用对；P2 决定先研究哪个商品范围；P3 规定“什么算一个问题”；P4 让人能快速找到原始证据；P5 比较不同查找方法；P6 才让模型按步骤调用工具；P7 检查目标运营人员是否真的愿意使用。后面每一步都在解决前一步留下的具体问题。
\end{beginnerbox}

\begin{beginnerbox}
这里的“证据”不是指系统已经证明某个商品真的存在质量缺陷，而是指：系统给出的判断能够返回到本地、冻结数据中的对应评论记录，供运营人员继续复核。评论信号是后续调查的输入，不是现实世界结论。
\end{beginnerbox}

\section{开始前需要知道的状态}

\begin{factbox}
仓库当前可验证状态为 P0--P6 \statuspassed；P7 为 \statusprogress。

合成浏览器流程和受控本地工程回归已有通过证据，但真人目标运营用户尚未完成结构化试用。因此，不能把这一阶段写成通过，也不能声称项目已经为企业降本增效。
\end{factbox}

阶段状态和放行条件以 \filelink{docs/master-plan.md} 为执行权威；当前状态摘要见 \filelink{docs/current-stage-status.md}；阶段学习、失败复盘与面试边界见 \filelink{docs/interview-casebook.md}。阶段报告位于 \filelink{docs/reports/}。这些都是仓库相对路径，不依赖某台电脑的目录结构。

\section{四类信息必须分开}

阅读本手册时，不要把“看起来合理”当成“已经证明”。正文使用以下四种信息类型。

\begin{factbox}
\textbf{事实：}已有仓库产物、测试、审计或阶段报告支持。例如，P6 的冻结案例工程门已经通过。这类表述仍然受数据快照、案例范围和配置版本限制。
\end{factbox}

\begin{inferencebox}
\textbf{推断：}从已有事实出发得到的合理解释。例如，评论证据工作台可能减少重复查找和整理评论的机械操作。它有理由成立，但在真人任务试用前还不是业务结果。
\end{inferencebox}

\begin{advicebox}
\textbf{建议：}下一步应做的产品或工程动作。例如，先让目标运营用户完成同一组手工任务和工作台任务，再比较完成时间与证据复核情况。建议不是已经完成的工作。
\end{advicebox}

\begin{boundarybox}
\textbf{待验证：}当前证据不足以确认的信息，包括真实需求频率、用户采纳率、付费意愿、实际节省时间、决策质量改善和生产环境可靠性。证据不足不等于这些结论一定错误，只表示现在不能声称它们成立。
\end{boundarybox}

\section{先理解 P0 与 P1：数据地基}

第一次阅读时，可以把 P0 和 P1 当成一个学习单元，连续回答两个问题：

\begin{enumerate}
\item P0：这批数据能不能放心用？
\item P1：之后的分析是不是一直在用同一版数据？
\end{enumerate}

举个运营例子：你今天发现某类商品的“漏水”反馈比较集中。下周重新运行时，如果换了一份文件，结果变化可能只是数据换了，并不代表商品真的变好或变坏。P0 先检查数据来源、字段和使用边界；P1 再给本次输入做版本记录。这样后面的结论才有机会比较和复盘。

\begin{beginnerbox}
P0/P1 的业务价值不是立刻生成销售结论，而是避免“做了很久才发现数据不能用”以及“每次运行都不知道自己用了哪批数据”。技术名词先不用背，等读到对应阶段时再看术语卡。
\end{beginnerbox}

\section{初学者术语表}

\begin{longtable}{>{\raggedright\arraybackslash}p{0.18\textwidth}>{\raggedright\arraybackslash}p{0.72\textwidth}}
\toprule
术语 & 通俗解释 \\
\midrule
\endfirsthead
\toprule
术语 & 通俗解释 \\
\midrule
\endhead
B 端 & 企业或团队内部使用的工具。本项目的主用户是产品或类目运营。 \\
C 端 & 普通消费者直接使用的产品。本项目不做消费者选购助手。 \\
SKU & 一个具体可区分的商品单元。手册中的“单品诊断”可以先理解为针对一个指定商品查找评论证据。 \\
VOC & Voice of Customer，即“客户之声”。在本项目中，它指从评论中整理出来的体验维度和问题信号。 \\
Agent & 能根据请求选择既定步骤和工具的软件流程。这里的 Agent 不能任意执行代码或查询，模型只提出受约束的计划和表达。读到 P6 再深入理解。 \\
确定性工具 & 输入相同时按照固定规则查询和计算的本地代码，例如商品解析、证据检索和引用校验。 \\
Verifier & 验证器。它不相信模型“说自己正确”，而是逐项核对商品范围、证据强度、引用和结果状态。 \\
Gate & 阶段闸门或放行条件。先把它理解成“这一关的检查清单”；只有规定的输入、产出和阻断指标都满足，阶段才可以标为通过。 \\
Lineage & 数据血缘。先把它理解成“结果是从哪里来的记录”：哪一批数据、哪套规则和哪份代码。 \\
Silver 标签 & 由冻结规则自动构造的期望结果，适合回归测试，但不是人工双审得到的 gold 标签。 \\
失败关闭 & 条件不满足时停止正常输出并返回明确错误或“证据不足”，而不是猜测一个看似完整的答案。 \\
\bottomrule
\end{longtable}

\section{推荐阅读顺序}

本手册采用一条按阶段推进的项目故事。第一次阅读请先读“项目为什么值得做”，然后依次读 P0--P7；P0/P1 作为“数据地基”连续理解，Agent 知识会在 P4--P6 的真实开发顺序中出现，而不是作为与证据系统分离的课程。

\begin{tabularx}{0.94\linewidth}{>{\raggedright\arraybackslash}p{0.18\textwidth}X}
\toprule
读者目标 & 建议路径 \\
\midrule
第一次接触项目 & 先读“为什么做”，再读 P0--P7，理解每一步在保护什么结论。 \\
完整学习工程 & 按 P0--P7 连续阅读；到 P6 时回看 P3--P5 的合同、工具与检索证据，理解模型为何不掌握最终事实。 \\
准备业务或产品面试 & 重点阅读用户、三个工作流、价值边界、阶段业务价值和真人验证缺口。 \\
准备技术面试 & 重点阅读数据血缘、确定性工具、Structured Output、Verifier、LangGraph、checkpoint 和评测失败复盘。 \\
准备简历 & 先确认当前状态和证据，再提炼“问题、行动、技术选择、可核实结果、边界”，不要只抄测试数字。 \\
\bottomrule
\end{tabularx}

\begin{advicebox}
基础较弱时，不需要先背框架名或哈希字符串。每读一个阶段，先回答三个问题：它要解决哪个运营问题？它留下了什么可以检查的结果？它还不能证明什么？能回答后，再学习对应技术会更稳固。
\end{advicebox}

\section{全书共同遵守的表达边界}

\begin{boundarybox}
本手册不得把自动化测试写成真人验收，不得把 silver 指标写成人工质量，不得把评论信号写成缺陷或因果事实，也不得把内部原型写成已上线企业 SaaS。公开材料只使用代码、合成夹具、聚合指标和经审查的结果；不展示原始评论、真实商品或证据标识、密钥和本地绝对路径。
\end{boundarybox}
